\section{Introduction}

\subsection{The aggregation problem}

Casualty figures in armed conflict are reported by sources with opposing incentives. Belligerents inflate enemy combatants killed and deflate civilians; advocacy organisations, opposing belligerents, and intergovernmental bodies have symmetric incentives in the other direction. Naive aggregation is therefore biased in a direction that depends on which side the analyst trusts more. The tempting conclusion is that civilian-to-combatant ratios are too politicised to analyse statistically.

We argue for a narrower position. Conflicting source aggregates fail in well-understood ways, but the parameters they purport to measure are jointly constrained by a small number of relations whose inputs are measured by different, non-aligned institutions. An analyst who treats every belligerent source as adversarial can still recover an informative bound by exploiting these constraints; what cannot be recovered without further assumptions is the position of the truth \emph{within} the bound. This separation---bounds from data, position from politics---is the organising principle of the paper. The framework is a \emph{diagnostic}, not an oracle: its purpose is to force any dispute about a casualty claim to name the specific independent measurement being contested, and to quantify how far that measurement would have to be wrong.

\paragraph{A worked example.} Suppose a population is $73\%$ women and children ($w=0.733$) and records of the dead show that $56\%$ of $70{,}000$ dead are women and children ($\omega=0.560$). If violence were blind to demographics, the dead would mirror the population, so the deficit of women and children among the dead is informative: some of the dead were combatants (all adult men, by assumption), and possibly civilian adult men were more exposed than their population share suggests. If civilian men and women die at \emph{equal} per-capita rates, arithmetic (Section~\ref{sec:bounds}) turns these three numbers into a combatant share of $q\approx 25\%$. If civilian adult men die at \emph{twice} the rate of women and children---the kind of differential found when the same method is calibrated on historical conflicts where the answer is known \citep{frost2026}---the same arithmetic gives $q\approx 6\%$. The interval $[{\sim}6\%,{\sim}25\%]$ is what the data alone deliver; a claimed $q$ of $30\%$ is outside it for \emph{every} admissible exposure ratio and can only be rescued by asserting that the census or the death records are wrong, by an amount the framework computes in standard-error units. That amount is the claim's \emph{contradiction radius}.

The closest precedents are \citet{spagat2009} on the contestation of war-death estimates, \citet{lumpricebanks2013} on multiple systems estimation under reporting failure, \citet{radford2023} on Bayesian inference of conflict losses with reporting biases, \citet{vesco2026} on probabilistic UCDP underreporting, and \citet{gomez2025} on uncertainty quantification for Gaza mortality. Methodologically we draw on partial identification \citep{manski2003,manski2007,molinari2020,klinetamer2023} and robust Bayesian analysis \citep{huber1964,huber1981,berger1986}. On Gaza specifically, complementary demographic decompositions of the identified dead have been developed by \citet{cockerill2024} and \citet{frost2026}, and we engage with both.

\subsection{Two deliverables}

\paragraph{Bounds.} Where can the civilian-to-combatant ratio possibly lie, given a set of independently measured demographic and manpower facts? The answer is a sharp identified set in the sense of Manski's partial-identification programme \citep{manski2003,tamer2010,romanoshaikh2010}. If the identified set is narrow, no source-aggregation rule can pull the truth outside it. The width is determined by the data, not by the analyst's politics.

\paragraph{Confidence coefficients.} How much can the analyst's political-bias prior on each source widen this set? Each source is Huber $\varepsilon$-contaminated \citep{huber1981,berger1986}: it is true with probability $1-\varepsilon$ and adversarial with probability $\varepsilon$. A bias-robust sensitivity band inflates the face-value interval by $\sum_i\varepsilon_i R_i$ to first order in $\varepsilon$, where $R_i$ is the diameter of the identified set when source $i$ is removed; conditional on a realised contamination of source $i$, the worst-case displacement is $R_i$ itself. Setting $\varepsilon_i\approx 1$ for belligerent claims and $\varepsilon_i\approx 0.05$ for census-grade sources gives an interval honest about which inputs are doing the work.

The two objects are complementary: the first is information that survives any reasonable politics; the second is the rate at which uncertainty grows as the analyst dials in distrust of specific sources.

\subsection{The under-fudgible mechanism}

The bounds are not vacuous because of demographic arithmetic. Consider any claimed decomposition of the dead into combatants and civilians. The claim must simultaneously respect (i) the demographic composition of the identified dead, (ii) the demographic composition of the pre-war population, (iii) the size of the armed group's pre-war stock, and (iv) the plausible range of the civilian adult-male exposure differential. These quantities are measured by \emph{different} institutions: a national or UN census produces (ii); a health ministry, the OHCHR, or a household survey produces (i); independent military-balance assessments produce (iii); and (iv) can be calibrated on historical conflicts where combatant status was recorded \citep{frost2026}. They are not equally manipulable, and---critically---they cannot all be moved at once without the moves being visible.

We call a relation between conflict statistics \emph{under-fudgible} when its inputs are reported by non-aligned institutions (or have uncontroversial physical bounds), so that adjusting any single input to rescue a disputed claim forces a measurable deviation in at least one independently measured quantity. This is a heuristic property of a measurement system, not a formal axiom; what is formal is its consequence, the contradiction radius of Section~\ref{sec:contradiction}, which quantifies exactly how large the forced deviation is for any specific claim. The trigger is \emph{joint inconsistency}, not the size of the claim: the same machinery rejects a claimed combatant share that is too high and one that is too low, against the same feasible set. This symmetry matters. The diagnostic responds to the demographic content of a claim, not to its political direction, and Section~\ref{sec:gaza} applies it in both directions.

\subsection{Contributions and roadmap}

The paper delivers, in order of importance:

\begin{enumerate}[leftmargin=*,label=(\roman*)]
\item \emph{A sharp identified set} for the combatant share $q=D_M/D$ under bounded differential exposure of civilian adult men, intersected with the manpower bound $D_M \le M$ (Section~\ref{sec:bounds}). The set is a closed-form interval computable from four numbers.
\item \emph{The contradiction radius} $\rho$: a feasibility test for any disputed claim, stated as a distance-to-feasibility optimisation in standard-error units (Section~\ref{sec:contradiction}). It converts ``we dispute your number'' into ``name which independent input is wrong, and by how many standard errors.''
\item \emph{A bias-robust sensitivity band} under Huber $\varepsilon$-contamination of each source: a patternwise worst-case envelope, a first-order band for the interval endpoints, and a quantile-displacement bound under an explicit posterior-weight condition (Section~\ref{sec:bias}).
\item \emph{A precision-weighted aggregation rule} for multiple demographic sources, with an explicit caveat for selection-biased sources (Section~\ref{sec:aggregation}).
\item \emph{A Gaza 2023--2026 application} (Section~\ref{sec:gaza}): the identified set, the contradiction radius of the public IDF combatant-death claim under every anchor and exposure assumption, a sensitivity grid over all contested inputs, and a comparison against every independent method known to us---which all reject the claim while disagreeing only within the identified set. A fully documented spatial Bayesian model (Appendix~\ref{app:spatial}) illustrates where a specific set of modelling assumptions lands \emph{within} the set.
\item \emph{An 81-conflict overview} (Section~\ref{sec:empirical}) mapping where the diagnostic can and cannot be applied sharply, with per-conflict tables in the supplement.
\end{enumerate}

Proofs are collected in Appendix~\ref{app:proofs}. The method does not assign moral responsibility, estimate indirect deaths, or replace forensic work.
